\documentclass{article}
\usepackage[T1]{fontenc}
\usepackage{amsmath}

\title{Assignment 7}

\begin{document}
\maketitle

\section*{Question 1}

\subsection*{(a)}

we are given:

\begin{itemize}
    \item intensity of the light: $I=1\,\mathrm{W\,m^{-2}}$
    \item binding energy: $E_{\mathrm{bind}}=1.8\,\mathrm{eV}$
    \item Bohr radius: $r=5.3\times10^{-11}\,\mathrm{m}$
\end{itemize}

we assume that the atom absorbs all the radiation that falls within its Bohr radius.

the area of the atom is approximately the area of a circle,

\[
A=\pi r^2
\]

substituting,

\[
A=\pi(5.3\times10^{-11})^2
\]
\[
A\approx8.82\times10^{-21}\,\mathrm{m^2}
\]

the intensity is power per unit area,

\[
I=\frac{P}{A}
\]

so,

\[
P=IA
\]

substituting,

\[
P=(1)(8.82\times10^{-21})
\]
\[
P=8.82\times10^{-21}\,\mathrm{W}
\]

the binding energy is given in eV, so we convert it to joules.

\[
1\,\mathrm{eV}=1.602\times10^{-19}\,\mathrm{J}
\]

therefore,

\[
E_{\mathrm{bind}}
=
1.8(1.602\times10^{-19})
\]
\[
E_{\mathrm{bind}}\approx2.884\times10^{-19}\,\mathrm{J}
\]

using

\[
E=Pt
\]
\[
t=\frac{E}{P}
\]

substituting,

\[
t=
\frac{2.884\times10^{-19}}
{8.82\times10^{-21}}
\]

\[
t\approx32.7\,\mathrm{s}
\]

hence,

\[
\boxed{t\approx33\,\mathrm{s}}
\]

so, it would take about 33 seconds to collect enough energy to eject an electron.

\subsection*{(b)}

the calculation in part (a) treats light as a continuous wave, so the electron slowly builds up energy until it has enough to escape.
 
in reality, light comes in packets called photons, each with energy $E=hf$. an electron absorbs one whole photon at a time. if one photon has enough energy ($hf\geq\phi$), the electron is ejected immediately.
 
making the light weaker only means fewer photons arrive, not that each photon has less energy. so emission still starts instantly.

\section*{Question 2}

\subsection*{(a)}

we have light travelling from a medium with refractive index

\[
n_1(\lambda)
\]

into a medium with

\[
n_2=1.
\]

using Snell's law,

\[
n_1\sin\theta_1=n_2\sin\theta_2
\]

since $n_2=1$,

\[
n_1\sin\theta_1=\sin\theta_2
\]

therefore,

\[
\theta_2=\arcsin(n_1\sin\theta_1)
\]

we want to find how much the refracted angle changes when the refractive index changes.

differentiate with respect to $n_1$,

\[
\boxed{
\frac{d\theta_2}{dn_1}
=
\frac{\sin\theta_1}
{\sqrt{1-n_1^2\sin^2\theta_1}}
}
\]

as $\theta_1$ increases, the denominator becomes smaller, so the change in $\theta_2$ becomes larger.

however, the light must still be transmitted into the second medium.

from Snell's law,

\[
\sin\theta_2=n_1\sin\theta_1
\]

and since $\sin\theta_2\leq1$, we have

\[
n_1\sin\theta_1\leq1
\]

the maximum possible value occurs when

\[
n_1\sin\theta_1=1
\]
\[
\sin\theta_1=\frac{1}{n_1}
\]
\[
\theta_1=\arcsin\left(\frac{1}{n_1}\right)
\]

this is the critical angle.

hence $\boxed{\theta_1=\theta_c}$ is the best choice of incident angle to maximise the dispersion.

at this angle, $\theta_2=90^\circ$

\subsection*{(b)}

we are given:

\[
n_{400}=1.8
\]

and

\[
n_{700}=1.7
\]

the incident angle should be the critical angle for the larger refractive index, since both wavelengths need to be transmitted.

from part (a),

\[
\theta_1 = \arcsin\left(\frac{1}{1.8}\right)
\]
\[
\theta_1\approx33.75^\circ
\]

\begin{itemize}
    \item for the $400\,\mathrm{nm}$ light,

\[
1.8\sin(33.75^\circ)=\sin\theta_{2,400}
\]

since $33.75^\circ$ is the critical angle,

\[
1.8\sin(33.75^\circ)=1
\]

so,

\[
\theta_{2,400}=90^\circ
\]

\item for the $700\,\mathrm{nm}$ light,

\[
1.7\sin(33.75^\circ)=\sin\theta_{2,700}
\]
\[
\sin\theta_{2,700}
=
1.7(0.5556)
\]

\[
\sin\theta_{2,700}\approx0.9444
\]
\[
\theta_{2,700}=\arcsin(0.9444)
\]
\[
\theta_{2,700}\approx70.8^\circ
\]
\end{itemize}

the angle between the two ends of the optical spectrum is therefore

\[
\Delta\theta
=
90^\circ-70.8^\circ
\]

\[
\Delta\theta\approx19.2^\circ
\]

\[
\boxed{\Delta\theta\approx19^\circ}
\]

\section*{Question 3}

the thin lens equation is

\[
\frac{1}{f}=\frac{1}{u}+\frac{1}{v}
\]

since the sun is very far away,

\[
u\rightarrow\infty
\]
\[
\frac{1}{u}\rightarrow0
\]

so,

\[
\frac{1}{f}=\frac{1}{v}
\]
\[
\boxed{v=f}
\]

so the image of the sun forms one focal length behind the lens.

the sun has an angular diameter of

\[
0.5^\circ
\]

so the angle from the centre of the sun to its edge is

\[
\theta=\frac{0.5^\circ}{2}=0.25^\circ
\]

consider the ray from the centre of the lens to the edge of the solar image.

using the right triangle,

\[
\tan\theta=\frac{r}{f}
\]
\[
r=f\tan\theta
\]

substituting $\theta=0.25^\circ$,

\[
r=f\tan(0.25^\circ)
\]

the diameter is

\[
D=2r
\]

therefore,

\[
D=2f\tan(0.25^\circ)
\]
\[
D=0.008726f
\]
\[
\boxed{D\approx8.73\times10^{-3}f}
\]

where $D$ and $f$ are in the same units.

\end{document}
